\section{Practice: The Cultivation of Habit and Household}
\label{sec:practice}

\noindent
The cycle turns to practice, the cultivation and the action through which a flourishing, cognised in its normative specification, is pursued. The task of this moment is the formation of the shared life such that the just eudaimonia of \xref{sec:cog-spec} is generated, and the governing constraint, carried from the prior paper, is that this formation is a cultivation and not a construction: the shared grammar of a relation cannot be legislated from a position of possession, but only tended into an openness within which it continues to be jointly reconfigured. Three frameworks are drawn upon at this moment, the virtue-theoretic, the evidential, and the phenomenological, and the moment is organised around the tension between cultivating a disposition and not legislating it.

\subsection{The virtue-theoretic component: flourishing as habituation}
\label{sec:prac-virtue}

The virtue-theoretic framework is the framework proper to the moment of practice, because it alone holds that flourishing is brought into being through the formation of disposition by habit. In the Aristotelian account, the excellences in accordance with which a life flourishes are acquired by habituation: one becomes just by doing just acts, temperate by doing temperate acts, until the disposition is settled into a \emph{hexis}, a stable state of character from which the corresponding action flows without struggle \citep{aristotle2009}. The practice of a shared flourishing is, on this account, the habituation of the two parties, severally and together, into the dispositions in accordance with which their shared life is good: the formation, by repeated practice, of the settled ways of attending, responding, and caring that constitute a good shared life.

The matter of this habituation is the everyday sensory and habitual fabric the prior paper identified as the matter of the field. A habit of a shared life is formed in the recurrence of concrete practices: a way of taking a meal together, a manner of greeting and parting, an ordering of the common rooms, a rhythm of attention to one another. These recurrences are the \emph{habitus} in the sociological sense, the durable dispositions, inculcated by a form of life, that generate the practices of that form of life and reproduce it \citep{bourdieu1990}; and they are the family manner, the customs and bearing transmitted across a shared life and, in time, across generations. The cultivation of a shared flourishing is conducted in this fabric: not in the formulation of principles for the relation, but in the formation, by recurrence, of the dispositions borne in the concrete texture of the days.

The mechanism by which a disposition forms from a recurrence is worth making explicit, because it governs what the cultivation of a shared life can and cannot do. A practice repeated becomes easier, then expected, then a part of how the shared life is, until at last it is performed without deliberation and felt as natural, the disposition having settled into the \emph{hexis} from which it flows unbidden. This is why the matter of a shared flourishing is the small and recurrent rather than the large and singular: it is not the grand resolution but the repeated practice that forms a disposition, not the declaration of how the two will live but the daily living of it that settles into the family manner they come to have. A couple does not become attentive to one another by resolving to be attentive; they become attentive by attending, in the concrete recurrences of their days, until attention is the disposition from which they act. The cultivation of a shared flourishing works, accordingly, on the recurrences: it tends the daily fabric such that the practices worth settling into dispositions are the ones that in fact recur, and it does so knowing that what recurs will become, in time, what the parties are. This is also why the corruption of a shared life is so often gradual and unremarked, a small unkindness or a small neglect, recurring, settling by the same mechanism into a disposition, until the relation has acquired a character no one resolved upon and each performs without deliberation. The formation of disposition by habituation is neutral as to the disposition formed, and the cultivation of a shared flourishing is the tending of the recurrences such that the dispositions formed are the excellences of \xref{sec:cog-virtue} and not their contraries.

\subsection{The constraint: cultivation without legislation}
\label{sec:prac-constraint}

The formation of a shared disposition stands under a constraint that distinguishes cultivation from its corruption. The prior paper established that to possess and modify the generative structure of a relation from a position claimed to be external to it is exploitation in the symbolic register, the appropriation of a jointly generated generativity to a single legislator, who thereby reduces the other from a co-author of the shared life to an object governed by it. Applied to the formation of habit, this yields the constraint that the cultivation of a shared disposition must guide the relation's own dynamics without legislating its productions: it must tend the conditions under which good habits form and reform themselves jointly, and it must not impose a fixed habit as a rule the other is required to obey.

\begin{claim}[Cultivation guides the formation of habit; legislation possesses it]
The cultivation of a shared habit is the tending of the conditions under which a disposition forms and continues to reform itself in the joint practice of both parties; it guides the relation's own habit-forming dynamics without fixing their product. The legislation of a habit is the imposition, by one party from a position claimed to be external, of a fixed disposition as a rule the other is required to keep; it possesses the relation's habit-forming dynamics and fixes their product. The two are distinguished as the cultivation of a field is distinguished from the construction of a grammar, and as guidance is distinguished from possession in the prior paper's account of symbolic justice. A habit that is cultivated remains jointly authored and open to joint reform, and the disposition it forms is the disposition of both; a habit that is legislated is the disposition the legislator has fixed, and its imposition reduces the other to its object. The formation of a shared disposition is just when it cultivates and unjust when it legislates, and the difference is not in the content of the habit but in whether the habit-forming dynamics is guided or possessed.
\end{claim}

\noindent
This is the central tension of the moment of practice, and it admits of no general procedure for its resolution, since a procedure for forming the right habits, applied from outside the joint practice, would itself be a legislation. The tension is managed, in the manner of the prior paper's practice, subtractively: by the cultivation of the conditions under which good habits arise rather than the imposition of the habits themselves, by the preservation of the other party's capacity to reform a shared habit rather than the fixing of it, and by the recurrent question whether a habit of the shared life is still jointly authored or has hardened into a rule one party keeps and the other imposes. A family manner is well cultivated when it remains the joint and revisable bearing of both, and ill cultivated when it sets into a regime one administers and the other obeys.

\subsection{An operational vocabulary for cultivation without legislation}
\label{sec:prac-vocabulary}

The constraint is stated as a principle, and a foundational treatment owes the subsequent papers an operational vocabulary in which the principle can be applied, since a principle that cannot be brought to bear on the recurring practice of a shared life is of little use to a praxis. This subsection supplies such a vocabulary, under the operational concession: the terms it introduces are operational instruments for the conduct of a practice, and not measurements of the unsymbolised value the practice serves.

The first instrument is a set of signals by which the drift from cultivation toward legislation may be recognised before it has hardened. The drift is gradual, by the mechanism of habituation (\xref{sec:prac-virtue}), and it is therefore most often unremarked; the signals are the marks by which it may be caught. A shared habit is drifting toward legislation when its revision by one party comes to be treated by the other as a transgression rather than as a contribution; when the justification offered for a practice ceases to be a reason the other could weigh and becomes an appeal to how things are done; when one party's deviation from a practice is met with correction rather than with inquiry; when the practices of the shared life are increasingly described by one party in the language of rule, expectation, and obligation rather than in the language of joint preference; and when the question whether a practice still suits both parties is met with the answer that it is the practice, as though its being established were a reason for its continuance. These signals are operational and defeasible, and no one of them is decisive; their accumulation is the operational sign that a habit is hardening from a joint bearing into an administered rule, and that the cultivation has begun to legislate.

The second instrument is a repertoire of subtractive operations, the operations by which a legislating practice is returned toward cultivation. The repertoire is subtractive because the constraint forbids the additive remedy, the imposition of a better rule, which would be a further legislation; what cultivation can do is remove the conditions under which the legislation operates rather than impose a corrected practice. The subtractive operations include the suspension of a hardened practice, its deliberate setting-aside for a season so that the parties may discover whether it is jointly wanted or merely established; the reopening of a foreclosed question, the deliberate return of a settled practice to the status of an open joint preference; the withdrawal of the language of rule, the deliberate redescription of a practice from obligation to preference; and the restoration of the other party's standing to revise, the deliberate cession, by the party who has come to administer a practice, of the authority to do so. These operations do not impose a just practice; they remove the conditions under which an unjust one is maintained, and return the relation's habit-forming dynamics to the joint authorship from which a just practice may, but need not, arise.

The third instrument is the preservation of betrayability, which the prior paper identified as the mark of a relation that has not foreclosed itself. A shared habit preserves betrayability when its continuance remains contingent on the continued joint willing of both parties, so that either could, in principle, decline it, and its continuance is therefore a renewed joint choice rather than a settled fact. A habit that has lost betrayability is one whose continuance no longer depends on being willed, that persists by its mere establishment and would persist against the will of a party who came to decline it; and such a habit, however benign its content, has become a legislation, since it binds independently of joint willing. The preservation of betrayability is the operational form of the cultivation-not-legislation constraint at the level of the individual habit: a practice is cultivated, and not legislated, to the extent that it remains betrayable, that its continuance remains a renewed joint choice and does not harden into a fact that binds without being willed. The maintenance of betrayability across the habits of a shared life is, accordingly, the central operational task of a just practice, and the signals and subtractive operations above are the instruments of its maintenance.

These three instruments, the signals of drift, the subtractive operations, and the preservation of betrayability, compose an operational vocabulary in which the cultivation-not-legislation constraint can be applied to the recurring practice of a shared life. They are offered to the subsequent papers as the operational form of a principle the present paper states, and they stand, like all the paper's operational apparatus, under the concession that they are instruments of a practice and not measurements of the value the practice serves.

\subsection{The evidential component: supported practice and its limit}
\label{sec:prac-evidence}

The evidential framework contributes to the moment of practice the resources of evidence-based practice: the body of controlled findings concerning which practices conduce, on average and in studied populations, to the maintenance and the deepening of a shared life. Such findings are a genuine resource. That certain practices of attention, of repair after conflict, of the just division of a shared life's labour, conduce on the evidence to its flourishing is a real gain for practice, that much of this sustaining labour is emotional labour, and is systematically under-recognised as labour, is argued by \citet{hochschild1983} and \citet{folbre2001}; and to disregard it on the ground that each relation is singular would be to discard a public and accumulable knowledge for a private intuition that is frequently mistaken.

The evidential component carries a limit that the operational concession has prepared. Evidence supports a practice by establishing that it produces a measured effect in a studied population; it speaks, therefore, in the register of the structured proxy, and of the average. It cannot establish what the cultivation of this shared life, in its singularity and its unsymbolised texture, requires, and it cannot reach the unsymbolised value in which the relation's good finally consists. The evidential component is thus a resource to be drawn upon and not a sufficient guide: it furnishes supported practices that the cultivation may employ, while the determination of what this shared life requires remains a matter of the joint and largely unsymbolised appraisal of \xref{sec:appraisal}, to which the evidence is an input and not a substitute.

\subsection{The phenomenological component: the cultivation of attention}
\label{sec:prac-phenom}

The phenomenological framework contributes to practice the cultivation of attention: the formation, as itself a habit, of a quality of attending to the other and to the shared life. The practice of a shared flourishing is not only the formation of dispositions to act but the formation of a disposition to attend, a settled readiness to perceive the other's state, the texture of the shared life, and the small alterations by which a relation rises or declines. This is the practical correlate of the basal appraisal of \xref{sec:appraisal}: since the relation is maintained by the continuous appraisal of its state, the cultivation of the quality of that appraisal, the formation of a fine and unbroken attention to the shared life, is among the principal practices of its flourishing. The cultivation of attention is, moreover, the practice least assimilable to legislation, since attention cannot be imposed as a rule but only formed as a disposition, and it is therefore the practice in which the cultivation-not-legislation constraint is most nearly automatic.

\subsection{The frameworks in concert}
\label{sec:prac-concert}

The three components compose the practice of a shared flourishing. The virtue-theoretic framework supplies its form, the formation of disposition by habituation in the concrete fabric of the shared life. The evidential framework supplies supported practices the cultivation may draw upon, within the limit that they are proxies and averages and not the singular requirement of this life. The phenomenological framework supplies the cultivation of attention, the disposition to appraise the shared life finely and without cease. And the whole stands under the constraint that the formation of disposition is a cultivation and not a legislation, a guidance of the relation's own habit-forming dynamics and not the possession of it. The practice so composed generates the value whose circulation the next moment of the cycle assesses; the cycle now turns to the moment of evaluation, at which the question is how a shared life so cultivated is to be assessed, and at what cost to the value the assessment cannot reach.
